\section{The marginal and exclusion of every Monge optimizer}
\label{sec:proof}

We now prove Theorem~\ref{thm:main}. Let the balls, maps, and
potential be as in Proposition~\ref{prop:certificate} and
Lemma~\ref{lem:common-ball}. Set
\[
 \calH=\{Y_1,W_1,Y_2,W_2\}.
\]

\paragraph{A smooth marginal.}
Choose a nonnegative function \(g\in C_c^\infty(B)\) with
\(\int_Bg^2\dd x=1\), and let
\(\nu=g^2\dd x\), regarded as a probability measure on \(\R^3\).
Define
\begin{equation}\label{eq:mu}
 \mu=\frac13\nu+\frac16\sum_{H\in\calH}H_\#\nu.
\end{equation}
These coefficients sum to one. For \(H\in\calH\), define
\[
 g_H(y)=
 \begin{cases}
  g(H^{-1}(y))\,|\det DH^{-1}(y)|^{1/2},&y\in H(B),\\
  0,&y\notin H(B).
 \end{cases}
\]
The formula is smooth on \(H(B)\). Its support is contained in
the compact set \(H(\supp g)\subset H(B)\), so extension by zero
is smooth on \(\R^3\). Thus \(g_H\in C_c^\infty(\R^3)\), and
change of variables gives \(H_\#\nu=g_H^2\dd y\).

It follows that \(\mu=\rho\dd x\), where
\[
 \rho=\frac13g^2+\frac16\sum_{H\in\calH}g_H^2.
\]
The functions here are nonnegative and have pairwise disjoint
supports, since their supports lie in the five disjoint
components. Consequently
\[
 \sqrt{\rho}=\frac1{\sqrt3}g
                +\frac1{\sqrt6}\sum_{H\in\calH}g_H.
\]
This proves both asserted smoothness properties, compact support,
and normalization. Also \(\mu(U)=1\) and
\(\mu|_B=(1/3)\nu\).

\paragraph{A finite optimal coupling.}
Sample \(x\) with law \(\nu\), choose \(k\) uniformly from
\(\{1,2\}\), independently of \(x\), and apply an independent
uniform permutation to the three entries of
\[
 (x,Y_k(x),W_k(x)).
\]
Let \(\pi_0\) be the resulting law.
In any fixed coordinate the central point has coefficient \(1/3\),
and each specified outer map has coefficient
\((1/2)(1/3)=1/6\). Hence every coordinate marginal is exactly
\(\mu\).

All sampled triples satisfy \eqref{eq:certificate} with equality.
The boundedness of \(u\) therefore gives
\begin{equation}\label{eq:optimal-value}
 \int c\dd\pi_0=3\int_Uu\dd\mu<\infty.
\end{equation}
Conversely, every \(\pi\in\Pi(\mu,\mu,\mu)\) has
\(\pi(U^3)=1\): each of the three events \(x_i\notin U\)
has probability zero. Integrating \eqref{eq:certificate} yields
\[
 \int c\dd\pi\ge 3\int_Uu\dd\mu.
\]
Thus \(\pi_0\) attains \(\E(\mu)\), and the latter is finite.
Only the inequality on \(U^3\) is used.

For any optimal \(\pi\), the nonnegative slack
\(c(x_1,x_2,x_3)-\sum_i u(x_i)\) on \(U^3\) has integral zero.
Indeed the cost is finite and the potential sum is bounded, so
subtraction and integration are legitimate.
It follows that \(\pi\) is concentrated on the equality triples
in \eqref{eq:contact}.

\paragraph{No optimal maps.}
Suppose that Borel maps \(T_2,T_3\), each preserving \(\mu\),
induce an optimal plan. The preceding equality condition holds
at \((x,T_2(x),T_3(x))\) for \(\mu\)-almost every \(x\), hence
for \(\nu\)-almost every \(x\in B\).

The first coordinate \(x\in B\subset U_0\) is the unique central
coordinate of an equality triple. The remaining coordinates lie
one each in \(U_k^+\) and \(U_k^-\) for some \(k\).
If their values in those components are \(y,z\), then
\(x=X_k(y)\) and \(z=Z_k(y)\).
Injectivity of \(X_k\) forces
\[
 y=Y_k(x),\qquad z=W_k(x).
\]
Allowing either order of these coordinates and either value of
\(k\), we obtain
\[
 T_2(x)\in\{Y_1(x),W_1(x),Y_2(x),W_2(x)\}
 \quad\text{for \(\nu\)-almost every }x\in B.
\]
This contradicts Lemma~\ref{lem:splitting} applied to
\eqref{eq:mu}. Thus no Monge plan attains the minimum, completing
the proof of Theorem~\ref{thm:main}.\qed

